\subsection{Related Work on RL for LLM Reasoning}
\label{sec:related-work-rl-reasoning}

The release of DeepSeek-R1 and DeepSeek-R1-Zero in January 2025 ignited a wave of research into reinforcement learning (RL) as a mechanism for eliciting and scaling reasoning capabilities in large language models. Within weeks of the original publication, numerous independent teams began reproducing, extending, and analyzing the R1 training paradigm---often at dramatically smaller model scales and compute budgets than the original. This section surveys the most influential of these concurrent efforts, which collectively establish the empirical foundations upon which the present work on GRPO scaling builds.

\subsubsection*{DeepScaleR: Iterative Context Lengthening for Efficient RL Scaling}

DeepScaleR~\cite{luo2025deepscaler}, developed by Luo, Tan, Wong, Cai, Shi, Tang, Roongta, Zhang, Li, Popa, and Stoica at UC Berkeley and Agentica, demonstrated that reinforcement learning can be scaled effectively even with remarkably small models. The central contribution is a \emph{1.5B parameter} model, DeepScaleR-1.5B, fine-tuned from DeepSeek-R1-Distill-Qwen-1.5B using GRPO with a novel \textbf{iterative context lengthening} strategy. Rather than training with a fixed, long context window from the outset---which is computationally expensive and can lead to unstable training dynamics---the method begins with shorter context windows (8K tokens) and progressively extends them to 16K and then 24K tokens throughout training. This curriculum-style approach enables the model to first learn efficient reasoning patterns before learning to reason over longer chains of thought.

The results were striking: DeepScaleR-1.5B achieved \textbf{43.1\% Pass@1 on AIME 2024}, representing a 14.3 percentage point improvement over its base model and matching the performance of OpenAI's o1-preview---a model orders of magnitude larger. On MATH-500, DeepScaleR-1.5B achieved competitive results that further validated the effectiveness of RL scaling at small model sizes. The project also curated a training dataset of approximately 40,000 mathematics problems drawn from AIME (1984--2023), AMC, Omni-MATH, and other competition sources. Critically, the entire training recipe---including dataset, code, hyperparameters, and training logs---was open-sourced, establishing a reproducible baseline that subsequent work (including the present study) has built upon. The iterative context lengthening strategy has since been adopted by multiple follow-up projects, including FastCuRL~\cite{song2025fastcurl} and DeepCoder~\cite{deepcoder2025}, confirming its generality beyond the original mathematical reasoning setting.

\subsubsection*{SimpleRL-Zoo: Systematic Investigation of Zero RL Training}

While DeepScaleR focused on scaling a single model family, SimpleRL-Zoo~\cite{zeng2025simplerl} by Zeng, Huang, Liu, Liu, He, Ma, and He undertook a broader empirical investigation of \emph{zero RL training}---the paradigm where reinforcement learning is applied directly to base models without any supervised fine-tuning, relying solely on rule-based rewards. Published at COLM 2025, this work systematically studied zero RL training across \textbf{10 diverse base models} spanning different families and scales, including LLaMA-3-8B, Mistral-7B/24B, DeepSeek-Math-7B, Qwen2.5-Math-7B, and all Qwen2.5 models from 0.5B to 32B.

The key findings challenged several assumptions in the community. First, the authors observed that most prior reproduction efforts had focused exclusively on the Qwen2.5 model series, which they found to be unrepresentative: Qwen2.5 base models already exhibited strong instruction-following and self-reflection abilities, likely due to pretraining data contamination or design choices, making them an unusually favorable starting point for zero RL training. Second, through careful ablation, the authors identified several critical design strategies: \textbf{(i)} adjusting the format reward to encourage structured reasoning without over-constraining generation, \textbf{(ii)} controlling query difficulty to maintain a meaningful learning signal throughout training, and \textbf{(iii)} careful entropy management to prevent premature convergence.

Perhaps most importantly, SimpleRL-Zoo's analysis of training dynamics revealed that different base models exhibit fundamentally distinct patterns during RL training. The commonly cited correlation between increased response length and emergence of cognitive behaviors such as self-verification (the so-called ``aha moment'') does not hold universally. Some models increased response length without developing genuine verification strategies, while others exhibited verification behaviors with only modest length increases. Notably, the authors observed the ``aha moment'' for the first time in small models outside the Qwen family, demonstrating that this phenomenon is not architecture-specific. These findings underscore the importance of base model quality---rather than model family or scale alone---as the primary determinant of successful RL training for reasoning.

\subsubsection*{TinyZero: Minimal Reproduction of R1-Zero Reasoning}

TinyZero~\cite{tinyzero2025}, developed by Pan, Zhang, Wang, Yuan, Peng, and Suhr at UC Berkeley, represents perhaps the most accessible entry point for understanding RL-driven reasoning emergence. The project reproduced the DeepSeek-R1-Zero training paradigm using Qwen2.5 base models at the \textbf{3B parameter} scale, built upon the veRL training framework, with the explicit goal of demonstrating that emergent reasoning behaviors can arise at minimal cost---under \$30 of compute.

The task domain was the \textbf{Countdown number game}, where players must combine a set of given numbers using basic arithmetic operations to reach a target value. Despite the apparent simplicity of this task, TinyZero demonstrated that through RL training with rule-based outcome rewards, the 3B base model \emph{spontaneously} developed sophisticated reasoning strategies including self-verification (checking intermediate results), systematic search (exploring multiple solution paths), and iterative revision (backtracking and correcting errors). These emergent behaviors were not explicitly trained or prompted---they arose purely from the optimization pressure of the outcome reward.

The ablation experiments yielded several practical insights for the field. Base model size proved critical: the 0.5B model could only guess solutions without learning to search or verify, while from 1.5B onward, models began developing search and self-verification strategies, with the 3B model exhibiting the most sophisticated reasoning. Both base and instruct models were viable starting points, with instruct models learning faster but converging to similar final performance. TinyZero's significance lies not in achieving state-of-the-art benchmark scores, but in demonstrating the \emph{minimal conditions} under which RL can induce genuine reasoning---providing important evidence that the phenomenon observed in DeepSeek-R1-Zero is robust and reproducible at small scale.

\subsubsection*{DeepCoder: Extending RL Reasoning to Code Generation}

DeepCoder~\cite{deepcoder2025}, a joint effort by the Agentica team and Together AI led by Luo, Tan, Huang, Patel, Ariyak, Wu, and collaborators, extended the RL scaling paradigm from mathematical reasoning to \textbf{competitive programming}. Built on a 14B parameter model fine-tuned from DeepSeek-R1-Distill-Qwen-14B, DeepCoder-14B-Preview achieved \textbf{60.6\% Pass@1 on LiveCodeBench} (covering problems from August 2024 to February 2025), representing an 8 percentage point improvement over the base distilled model and matching the performance of OpenAI's o3-mini (Low) and o1.

The technical approach employed an enhanced variant of GRPO, dubbed \textbf{GRPO+}, which incorporated several stabilizing modifications inspired by DAPO~\cite{yu2025dapo}: offline difficulty filtering to maintain an appropriate challenge level without the runtime overhead of online rejection sampling, removal of the KL penalty to allow greater policy divergence and exploration, and a token-level loss normalization scheme. The training adopted the iterative context lengthening strategy pioneered by DeepScaleR, scaling from 16K to 32K tokens, with inference-time extension to 64K. The training dataset comprised approximately 24,000 unique problem-test pairs from Taco-Verified, PrimeIntellect SYNTHETIC-1, and LiveCodeBench (May 2023--July 2024), with test-case execution providing the verifiable reward signal.

DeepCoder's significance extends beyond its benchmark results. It demonstrated that the RL scaling recipe developed for mathematical reasoning transfers effectively to the code domain, where reward signals are naturally verifiable through program execution. The project highlighted a key challenge specific to code RL: constructing high-quality datasets with reliable, comprehensive test suites is substantially more difficult than curating mathematical problem-answer pairs, as incomplete test coverage can lead to reward hacking where models produce solutions that pass available tests without being truly correct.

\subsubsection*{Open-Reasoner-Zero and Other Open Reproductions}

Open-Reasoner-Zero~\cite{hu2025openreasonerzero}, introduced by Hu, Zhang, Han, Jiang, Zhang, and Shum, distinguished itself as the first fully open-source implementation of large-scale reasoning-oriented RL training on base models with an explicit focus on \textbf{scalability, simplicity, and accessibility}. Using Qwen2.5-32B-Base---the same foundation as DeepSeek-R1-Zero-Qwen-32B---the project demonstrated that a minimalist approach suffices: vanilla PPO with Generalized Advantage Estimation ($\lambda=1$, $\gamma=1$) and straightforward rule-based rewards, \emph{without any KL regularization}, successfully scaled both benchmark performance and response length, replicating the scaling phenomenon observed in DeepSeek-R1-Zero.

The results were notable for their efficiency: Open-Reasoner-Zero achieved \textbf{superior performance} to DeepSeek-R1-Zero-Qwen-32B across AIME 2024, MATH-500, and GPQA Diamond, while requiring only \textbf{1/10 of the training steps} compared to the DeepSeek-R1-Zero pipeline. The analysis contributed an important mechanistic insight: the learned critic in the PPO framework effectively identifies and devalues repetitive response patterns, yielding more robust advantage estimations and enhancing training stability. This finding suggests a concrete mechanism by which PPO-based approaches may offer advantages over GRPO in certain settings, as the critic provides an additional channel for the model to learn which reasoning patterns are productive versus degenerate.

Beyond Open-Reasoner-Zero, several other open reproduction efforts have contributed to the collective understanding. The ``Understanding R1-Zero-Like Training'' study by Liu et al.~\cite{liu2025understanding} provided a critical perspective, revealing that GRPO exhibits an optimization bias that artificially increases response length (especially for incorrect outputs) during training, and proposed Dr.\ GRPO as an unbiased correction. Sky-T1 demonstrated that o1-preview-level performance could be achieved with minimal resources, while Logic-RL explored RL training on logical reasoning tasks beyond mathematics, finding that structured logical domains can also benefit from the R1-style training paradigm.

\subsubsection*{Scaling Trends in RL Post-Training}

The proliferation of RL-for-reasoning projects in early-to-mid 2025 has enabled the first systematic studies of scaling behavior specific to RL post-training. Tan et al.~\cite{tan2025scaling} conducted a comprehensive empirical investigation across the full Qwen2.5 dense model series (0.5B to 72B parameters), characterizing how model scale, data volume, and computational budget interact during RL-based post-training for mathematical reasoning. Their analysis yielded four key findings with direct implications for practitioners:

\begin{enumerate}
    \item \textbf{Compute efficiency scales with model size:} Under a fixed computational budget, larger models trained for fewer RL steps consistently outperform smaller models trained for more steps, mirroring the Chinchilla-style findings from pretraining but now established for RL post-training.
    \item \textbf{Power-law predictability:} The relationship between test loss, compute, and data during RL post-training can be modeled by a predictive power-law, which holds robustly across both base and instruction-tuned starting points.
    \item \textbf{Diminishing returns at extreme scale:} While larger models exhibit higher learning efficiency, the analytical learning efficiency term $k(N)$ reveals a latent saturation trend as model size continues to increase, suggesting that simply scaling model parameters yields diminishing marginal returns for RL post-training.
    \item \textbf{Data reuse is effective:} In data-constrained regimes, repeated reuse of high-quality training data proves highly effective, as final RL performance is primarily governed by the total number of optimization steps rather than the uniqueness of training samples.
\end{enumerate}

These scaling laws complement the findings from individual projects reviewed above. DeepScaleR and TinyZero demonstrated that RL can be effective at surprisingly small scales (1.5B and 3B parameters, respectively), while Open-Reasoner-Zero showed that efficiency gains are possible even at larger scales (32B) through algorithmic simplicity. SimpleRL-Zoo emphasized that the \emph{quality} of the base model is often more important than its \emph{size}, and DeepCoder confirmed that these scaling principles generalize from mathematics to code. Collectively, these works establish that RL post-training for reasoning follows predictable scaling dynamics, but that the relationship between scale and capability is mediated by factors including base model pretraining quality, task domain, reward design, and training curriculum---factors that the present work on GRPO scaling seeks to further disentangle.
